\section{Discussion and Threats to Validity}
\label{sec:discussion}

\textbf{Measurement is the stronger contribution.} The clearest result is that an unprivileged canary can stand in for an un-instrumented foreground ($\rho=0.96$, AUC $\geq0.965$, 1.2\,\% false alarms), and that its components name the contended resource in 86--99.6\,\% of stalled periods. Windows exposes nothing like this to user mode today. A PSI-like signal is useful on its own: for diagnosis, for telemetry, and as input to any policy, including one that only advises the user.

\textbf{Static priority is a strong baseline.} On this host, lowering background priority recovers most of the attainable foreground latency for CPU, memory-bandwidth, and mixed contention, at modest cost. SmartSwap~v2's advantages over it are qualitative rather than in p95. It acts only when WSI shows harm. It chooses the I/O hint only when storage is the cause. It verifies and reverts every change. It records evidence that the change helped. Practitioners should try priority demotion first; our ladder does exactly that.

\textbf{A benefit check needs the right target.} Under memory-bandwidth contention the cap lowered WSI by enough to pass the benefit check but did not improve the foreground further. Because the canary is more sensitive than the foreground to residual contention, a WSI-only SLO can over-throttle. Two remedies follow from the data. One is a cost-aware benefit test that keeps a cap only if the WSI reduction per unit of background throughput lost exceeds a bound; SmartSwap already measures the group's CPU time, so this is computable. The other is a resource-specific SLO.

\textbf{Statistical power.} Ten blocks give large, consistent effects with bootstrap intervals that exclude no effect, but a 37-test Holm family cannot reach $p<0.05$ with $n=10$. A confirmatory study should pre-register a small primary family or use about 20 blocks.

\textbf{External validity.} All results come from one hybrid laptop on AC power in the \emph{Balanced} plan, with a busy but uncontrolled desktop. The foreground is a synthetic probe rather than a real application, and the background workloads are synthetic. Storage interference was weak on a fast NVMe drive and may be stronger on SATA SSDs or hard disks. Controller parameters were chosen in pilot runs on the same host; they were frozen before the study, but they were not validated on unseen hosts. Replication on homogeneous-core and AMD hosts, real foreground applications (input-to-display latency), and real background agents (sync, indexing, builds) are the most important next steps.

\textbf{Scope of control.} SmartSwap governs only processes it starts or is explicitly given. Choosing \emph{which} third-party processes to govern raises policy questions (consent, fairness, security software) that we leave open.

\section{Conclusion}
SmartSwap shows that stock Windows can measure the time interactive work loses to background interference from user mode, name the contended resource, and apply verified, reversible controls, with no privileges. On a hybrid laptop, the Windows Stall Index tracked an un-instrumented foreground's tail latency closely, and attribution was accurate. A least-cost ladder matched the strongest static baseline while acting only on evidence; removing the ladder made tail latency up to $4.8\times$ worse. We also documented how a first evaluation produced convincing but invalid results, and how the second one guards against each failure.
